\chapter{Pemphigus Vulgaris}
\index{Pemphigus Vulgaris}
\label{sec:Pemphigus Vulgaris}

\section{Overview}
Pemphigus vulgaris (PV) is a potentially life-threatening autoimmune blistering disease characterized by IgG autoantibodies against desmoglein 3 (and sometimes desmoglein 1), components of desmosomes that maintain keratinocyte adhesion.
\section{Causes}
This condition happens because of autoantibody-mediated loss of keratinocyte adhesion leading to intraepidermal blisters (suprabasal acantholysis).     
\section{Risk Factors}

\begin{itemize}[noitemsep]
  \item Genetic predisposition and family history
  \item Advancing age
  \item Lifestyle factors (diet, exercise, smoking, alcohol)
  \item Environmental exposures
  \item Underlying medical conditions
  \item Medication use
\end{itemize}
\section{Signs and Symptoms}

\subsection{Common symptoms}
\begin{itemize}[noitemsep]
  \item You may experience mucosal involvement (oral erosions) in 50--70\% of cases, often preceding cutaneous lesions
  \item Skin blisters are flaccid, easily ruptured (positive Nikolsky sign), and heal without scarring.
  \item Pain or discomfort in the affected area
  \end{itemize}

\subsection{Emergency warning signs}
\begin{itemize}[noitemsep]
  \item Severe or worsening symptoms of pemphigus vulgaris require immediate medical evaluation
\end{itemize}
\section{Progression and Stages}
The condition may progress through different stages of severity over time. Early stages often involve mild or subtle symptoms that may be overlooked. Moderate stages involve more noticeable symptoms affecting daily function. Advanced stages may involve significant impairment and complications.
\section{Complications}
\begin{itemize}[noitemsep]
  \item Disease progression leading to worsening symptoms
  \item Secondary infections or injuries
  \item Side effects of treatment
  \item Reduced quality of life
  \item Emotional and psychological impact
\end{itemize}
\section{Diagnosis}
\begin{itemize}[noitemsep]
  \item To make the diagnosis, doctors look for skin biopsy showing suprabasal acantholysis, direct immunofluorescence showing intercellular IgG and C3 deposition ("chicken-wire" pattern), and detection of anti-desmoglein antibodies by ELISA. Management requires immunosuppression: systemic corticosteroids combined with steroid-sparing agents (azathioprine, mycophenolate, rituximab)
  \item Rituximab (anti-CD20) has become a The first-choice treatment, often achieving drug-free remission.
  \item Comprehensive medical history and physical examination
  \item Laboratory tests to assess organ function
  \item Imaging studies to visualize affected structures
  \item Specialized testing depending on the affected system
  \item Consultation with relevant specialists
\end{itemize}
\section{Prevention}
\begin{itemize}[noitemsep]
  \item Maintain a healthy lifestyle with balanced diet and exercise
  \item Avoid known risk factors
  \item Attend regular medical check-ups for early detection
  \item Follow recommended screening guidelines
\end{itemize}
\section{Treatment Options}
Medications to manage symptoms and address underlying mechanisms Lifestyle modifications and self-care measures Physical therapy and rehabilitation Surgical intervention when indicated Regular monitoring and follow-up care Patient education and support services

\begin{itemize}[noitemsep]
  \item Medications to manage symptoms and address underlying mechanisms
  \item Lifestyle modifications and self-care measures
  \item Physical therapy and rehabilitation
  \item Surgical intervention when indicated
  \item Regular monitoring and follow-up care
\end{itemize}
\section{Living with It}
\begin{itemize}[noitemsep]
  \item Follow your treatment plan as prescribed
  \item Maintain regular follow-up with your healthcare provider
  \item Make necessary lifestyle adjustments
  \item Seek support from family, friends, or support groups
  \item Stay informed about your condition
\end{itemize}
\section{Outlook}
The outlook has improved significantly with modern immunosuppressive therapy. The outlook varies depending on the specific condition, severity at diagnosis, and response to treatment. Early intervention generally leads to better outcomes. Many conditions can be effectively managed with appropriate treatment. Regular follow-up care is important for monitoring and treatment adjustment.
